\subsection{Relatively bounded linear forms}

Let E be a directed ordered vector space. Let Q be the set of positive linear forms on E; it is a subset of the algebraic dual \(E^{*}\) of E (the space of all linear forms on E). It is immediate that \(Q + Q \subset Q\) and \(\lambda Q \subset Q\) for every scalar \(\lambda > 0\) (in other words, Q is a convex cone in \(E^{*}\) ). Moreover, \(Q \cap (-Q) = \{0\}\) , because if L and -L are both positive linear forms, then \(L(x) \geqslant 0\) and \(L(x) \leqslant 0\) for all \(x \geqslant 0\) , whence \(L(x) = 0\) for all \(x \geqslant 0\) and therefore L = 0 (No.~1, Prop. 3). The set Q thus defines on \(E^{*}\) an order relation \(L \leqslant M\) , equivalent to «M - L is a positive linear form on E», or again to «for every \(x \geqslant 0\) , \(L(x) \leqslant M(x)\) »; the elements \(\geqslant 0\) of \(E^{*}\) for this order structure are the positive linear forms (which justifies the terminology introduced). Let \(\Omega\) be the linear subspace of \(E^{*}\) generated by Q, that is, the set of linear forms on E that are differences of two positive linear forms; we are going to give another characterization of the elements of \(\Omega\) when E is a Riesz space.

DEFINITION 2. --- Given a Riesz space E, a linear form L on E is said to be relatively bounded if, for every \(x \geqslant 0\) in E, L is bounded on the set of \(y \in E\) such that \(|y| \leqslant x\) .

THEOREM 1. --- \(1^{\circ}\) In order that a linear form L on a Riesz space E be relatively bounded, it is necessary and sufficient that it be the difference of two positive linear forms.

\(2^{\circ}\) The ordered vector space \(\Omega\) of relatively bounded linear forms on E is a Riesz space that is fully lattice-ordered.

If \(L = U - V\) , where \(U\) and \(V\) are positive linear forms on \(\mathbf{E}\) , the relation \(-x \leqslant y \leqslant x\) implies

\[
- U (x) \leqslant U (y) \leqslant U (x) \quad \text { and } \quad - V (x) \leqslant V (y) \leqslant V (x),
\]

whence it is immediate that \(|L(y)| \leqslant U(x) + V(x)\) ; thus, L is relatively bounded. Suppose, conversely, that L is relatively bounded; it all comes down to proving that there exists a positive linear form N such that \(N(x) \geqslant L(x)\) for all \(x \geqslant 0\) , because N - L will then be a positive linear form.

Now, if a positive linear form N has this property then, for every \(x \geqslant 0\) and for \(0 \leqslant y \leqslant x\) , we have \(N(x) \geqslant N(y) \geqslant L(y)\) , therefore \(N(x) \geqslant \sup_{0 \leqslant y \leqslant x} L(y)\) ; if we prove that the real-valued function

\[
x \mapsto M (x) = \sup _ {0 \leqslant y \leqslant x} L (y),
\]

defined on the set P of elements \(\geqslant0\) of E, may be extended to a positive linear form on E (to be denoted M as well), we will have demonstrated the first part of the theorem and will have proved, moreover, that M is the supremum of 0 and L in \(\Omega\) . Since \(M(x)\geqslant0\) on P, it all comes down to proving that

\[
M (x + x ^ {\prime}) = M (x) + M (x ^ {\prime})
\]

for every pair of elements \(x \geqslant 0\) , \(x' \geqslant 0\) of E (No.~1, Prop. 3). By definition,

\[
\begin{array}{c} M (x) + M (x ^ {\prime}) = \sup _ {0 \leqslant y \leqslant x} L (y) + \sup _ {0 \leqslant y ^ {\prime} \leqslant x ^ {\prime}} L (y ^ {\prime}) \\ = \sup _ {0 \leqslant y \leqslant x, 0 \leqslant y ^ {\prime} \leqslant x ^ {\prime}} L (y + y ^ {\prime}) \leqslant M (x + x ^ {\prime}). \end{array}
\]

On the other hand, for every \(z\) such that \(0 \leqslant z \leqslant x + x'\) , we have \(x + x' = z + u\) with \(u \geqslant 0\) ; by the decomposition lemma (§1, No.~1), there exist two elements \(y, y'\) such that \(0 \leqslant y \leqslant x\) , \(0 \leqslant y' \leqslant x'\) and such that \(z = y + y'\) , \(u = (x - y) + (x' - y')\) ; then

\[
L (z) = L (y) + L \left(y ^ {\prime}\right) \leqslant M (x) + M \left(x ^ {\prime}\right),
\]

therefore \(M(x + x') = \sup_{0 \leqslant z \leqslant x + x'} L(z) \leqslant M(x) + M(x')\) , which completes the proof of the first part of the theorem. Moreover, we have thus shown that \(\Omega\) is a Riesz space and that, for every relatively bounded linear form L on E and for every \(x \geqslant 0\) ,

\[
L ^ {+} (x) = \sup _ {0 \leqslant y \leqslant x} L (y).\tag{1}
\]

It remains to see that \(\Omega\) is fully lattice-ordered; for this, it suffices to show that any set H of positive linear forms, bounded above and directed for the relation \(\leqslant\) , has a supremum in \(\Omega\) .

More generally, we have the following lemma:

Lemma. --- Let E be a directed ordered vector space, \(E^{*}\) its dual, ordered by taking as positive elements the positive linear forms. Let \((u_{\alpha})\) be an increasing directed family of elements of \(E^{*}\) . If, for every \(x \geqslant 0\) in E, \(\sup u_{\alpha}(x) < +\infty\) , then the family \((u_{\alpha})\) has a supremum u in \(E^{*}\) and, for all \(x \geqslant 0\) in E,

\[
u (x) = \sup _ {\alpha} u _ {\alpha} (x).\tag{2}
\]

In the set P of all \(x \geqslant 0\) in E, define the mapping u by the formula (2); it is immediate that \(u(\lambda x) = \lambda u(x)\) for all \(\lambda \geqslant 0\) and \(x \in P\) ; to prove the lemma it therefore suffices, by Prop. 2 of No.~1, to show that

\[
u (x + y) = u (x) + u (y)
\]

for \(x, y\) in P. But this is immediate on observing that \(u(x) = \lim u_{\alpha}(x)\) with respect to the directed set of indices (monotone limit theorem).

From the formula (1), one deduces immediately that if L and M are two relatively bounded linear forms on E then, for every \(x \geqslant 0\) ,

\[
\left\{ \begin{array}{l} \sup (L, M) (x) = \sup _ {y \geqslant 0,   z \geqslant 0,   y + z = x} \big (L (y) + M (z) \big) \\ \inf (L, M) (x) = \inf _ {y \geqslant 0,   z \geqslant 0,   y + z = x} \big (L (y) + M (z) \big). \end{array} \right.\tag{3}
\]

In particular if, in the first of these formulas, \(M\) is replaced by \(-L\) , we get

\[
| L | (x) = \sup _ {y \geqslant 0, z \geqslant 0, y + z = x} L (y - z).
\]

Now, if \(x = y + z\) , \(y \geqslant 0\) and \(z \geqslant 0\) , then \(-x \leqslant y - z \leqslant x\) ; conversely, the relation \(|u| \leqslant x\) implies \(L(u) \leqslant |L|(|u|) \leqslant |L|(x)\) . From this we deduce the formula

\[
| L | (x) = \sup _ {| y | \leqslant x} L (y) \quad \text { for } x \geqslant 0,\tag{4}
\]

whence, in particular,

\[
| L (x) | \leqslant | L | (| x |)\tag{5}
\]

for all \(x \in \mathbf{E}\) .

PROPOSITION 4. --- In order that two positive linear forms L, M on a Riesz space E be alien to each other in the space \(\Omega\) , it is necessary and sufficient that, for every number \(\varepsilon > 0\) and every \(x \geqslant 0\) in E, there exist two elements \(y \geqslant 0\) , \(z \geqslant 0\) of E such that \(x = y + z\) and \(L(y) + M(z) \leqslant \varepsilon\) . Indeed, by the second formula of (3), this condition expresses that \(\inf(L, M) = 0\) .

PROPOSITION 5. --- Let L be a positive linear form on a Riesz space E. In order that a positive linear form M on E belong to the band generated by L in \(\Omega\) , it is necessary and sufficient that, for every \(x \geqslant 0\) in E and every number \(\varepsilon > 0\) , there exist a number \(\delta > 0\) such that the relations \(0 \leqslant y \leqslant x\) and \(L(y) \leqslant \delta\) imply \(M(y) \leqslant \varepsilon\) .

Let us first show that the condition is necessary. If \(M \geqslant 0\) belongs to the band generated by L in \(\Omega\) , then (§1, No.~5, Cor. of Prop. 6)

\[
M = \sup _ {n} \left(\inf (n L, M)\right).
\]

If one sets

\[
U _ {n} = M - \inf (n L, M),
\]

\(U_{n}\) is thus a positive linear form on E and \(\inf_{n} U_{n} = 0\) in \(\Omega\) ; consequently (Lemma) \(U_{n}(x)\) tends to 0 as n tends to infinity, and there exists an n such that \(U_{n}(x) \leqslant \varepsilon / 2\) . Fixing such an n, we have \(U_{n}(y) \leqslant \varepsilon / 2\) for all y such that \(0 \leqslant y \leqslant x\) , thus the relation \(0 \leqslant y \leqslant x\) implies

\[
M (y) \leqslant \frac {\varepsilon}{2} + \inf (n L, M) (y) \leqslant \frac {\varepsilon}{2} + n L (y);
\]

if \(y\) is such that \(L(y) \leqslant \varepsilon / 2n\) it follows that \(M(y) \leqslant \varepsilon\) , which establishes our assertion.

Let us now show that the condition is sufficient. For every positive linear form \(M\) on \(\mathrm{E}\) , one can write \(M = U + V\) , where \(U\) belongs to the band generated by \(L\) in \(\Omega\) and where \(V\) is alien to \(L\) , \(U\) and \(V\) being positive (§1, No.~5, Th. 1). If \(M\) satisfies the condition of the statement then so does \(V = M - U\) , since \(0 \leqslant V \leqslant M\) . From this we will deduce that \(V = 0\) . For every \(x \geqslant 0\) in \(\mathrm{E}\) and every number \(\eta > 0\) , there exist two elements \(y \geqslant 0\) , \(z \geqslant 0\) of \(\mathrm{E}\) such that \(x = y + z\) and \(L(y) + V(z) \leqslant \eta\) (Prop. 4); given an arbitrary number \(\varepsilon > 0\) , choose \(\eta \leqslant \varepsilon\) so that the relations \(0 \leqslant u \leqslant x\) and \(L(u) \leqslant \eta\) imply \(V(u) \leqslant \varepsilon\) ; with \(y\) and \(z\) then determined as above, we have \(L(y) \leqslant \eta\) , therefore \(V(y) \leqslant \varepsilon\) and so

\[
V (x) = V (y) + V (z) \leqslant \varepsilon + \eta \leqslant 2 \varepsilon ;
\]

since \(\varepsilon\) is arbitrary, we have \(V(x) = 0\) for every \(x \geqslant 0\) , that is, \(V = 0\) .

Example. --- Let E be a Riesz space equipped with a locally convex topology compatible with its ordered vector space structure (TVS, II, §2, No.~7). Let \(E'\) be the topological dual of E, and suppose in addition that the cone P of elements \(\geqslant 0\) of E is complete for the weakened topology \(\sigma(\mathrm{E},\mathrm{E}')\) . Then every continuous linear form \(x' \in E'\) is relatively bounded, for one knows (TVS, II, §6, No.~8, Cor. 2 of Prop. 11) that under these conditions, for every \(x \geqslant 0\) in E the set of \(y \in E\) such that \(|y| \leqslant x\) is compact for \(\sigma(\mathrm{E},\mathrm{E}')\) . From this we deduce that E is then fully lattice-ordered; for (§1, No.~3, Prop. 2), it suffices to show that for every set \(H \subset E\) that is bounded above and directed for \(\leqslant\) , the section filter F of H is convergent in E for the topology \(\sigma(\mathrm{E},\mathrm{E}')\) (the latter being compatible with the ordered vector space structure of E). By translation, we can suppose that \(H \subset P\) , and it then suffices to show that F is a Cauchy filter for \(\sigma(\mathrm{E},\mathrm{E}')\) , or again that every continuous linear form \(x' \in E'\) has a limit with respect to F. But this follows at once from the monotone limit theorem when \(x'\) is a positive linear form, and since every linear form \(x' \in E'\) is the difference of two positive linear forms (Th. 1) our assertion is proved.
% semantic-fragments: legacy-input-seam
\relax{}
